\documentclass[12pt]{article}
\usepackage[utf8]{inputenc}
\usepackage[T1]{fontenc}
\usepackage{hyperref}
\title{Olùkùmi Language Resources for K-12 Curriculum Development}
\author{Compiled by Claude Sonnet 5 (claude-sonnet-5) for Andr\'as Kornai}
\date{}
\begin{document}
\maketitle

\section{Introduction}
Olùkùmi (ISO 639-3: ulb) is a Yoruboid language spoken in seven towns of Delta State, Nigeria: Ugbodu, Ukwu-Nzu, Ubulubu, Ugboba, Idumogo, Ogodor and Anioma. Sources describe it as endangered. The Ugbodu variety is closer to Yoruba and the Ukwunzu variety closer to Igbo.

\section{Primary Educational Resources}
No primers or school textbooks were found in the searches.

\section{Linguistic and Reference Materials}
\begin{itemize}
\item Olùkùmi Living Dictionary, which grew from a 2012 talking dictionary built with the Living Tongues Institute. It reportedly has over 1,600 entries and is multilingual with English, Yoruba and Igbo: \url{https://livingdictionaries.app/olukumi}
\item About page of the Living Dictionary: \url{https://livingdictionaries.app/olukumi/about}
\item Bolanle Arokoyo of the University of Ilorin is reported to have published a print dictionary in 2017. I did not find a link to it.
\item A Lexicostatistics Comparison of Yoruba, Igbo and Olukumi Dialects: \url{https://uilspace.unilorin.edu.ng/items/a7890ff3-de86-4cf9-b5ad-f0ee2f504d52}
\item Olukumi: a dialect of Yoruba in diaspora (University of Ibadan repository): \url{https://repository.ui.edu.ng/items/66b432f8-6419-4583-8e2e-7dc3952817cb}
\end{itemize}

\section{Religious and Community Resources}
None found.

\section{Digital and Audio Resources}
The Living Dictionary above is the main digital resource. It began as a talking dictionary, so it probably has audio, though I did not confirm this.

\section{Curriculum Development Context}
The language has a dictionary and some academic study but no found school materials. The online dictionary and Arokoyo's work could be a base for vocabulary and reading materials. The two varieties differ, so a standard form would need to be chosen with the community.
\end{document}